\chapter{集合论}
\intro{集合论是现代数学的基石，几乎所有数学对象都可以用集合来定义。本章讨论朴素集合论的核心概念，包括集合的基本定义、运算、关系以及容斥原理，为后续的关系、函数和图论奠定基础。}

\section{集合的基本概念}

\begin{mydefinition}{集合与元素}
\textbf{集合}是由确定对象构成的整体，构成集合的对象称为该集合的\textbf{元素}。用 $a \in A$ 表示 $a$ 是 $A$ 的元素，$a \notin A$ 表示 $a$ 不是 $A$ 的元素。
\end{mydefinition}

集合的基本特征：
\begin{enumerate}
\item \textbf{确定性}：对任意对象 $x$，$x \in A$ 和 $x \notin A$ 二者必居其一。
\item \textbf{无序性}：$\{1, 2, 3\} = \{3, 1, 2\}$。
\item \textbf{互异性}：$\{1, 1, 2, 2, 3\} = \{1, 2, 3\}$。
\end{enumerate}

\subsection{集合的表示法}

\textbf{列举法}：$A = \{1, 2, 3, 4, 5\}$；$\N = \{0, 1, 2, 3, \dots\}$。

\textbf{描述法}：$A = \{x \mid P(x)\}$，读作"满足 $P(x)$ 的所有 $x$ 构成的集合"。
例如 $A = \{x \mid x \in \R, x > 2\}$（大于 2 的所有实数）。

\textbf{文氏图}：用平面上的封闭曲线直观表示集合及其关系。用矩形表示全集 $U$，内部用圆表示各集合。

\subsection{特殊集合}

\begin{itemize}
\item \textbf{空集} $\varnothing = \{\}$：不含任何元素的集合。$|\varnothing| = 0$，对任意集合 $A$ 有 $\varnothing \subseteq A$。
\item \textbf{单元素集} $\{a\}$：仅含一个元素的集合。注意 $a \in \{a\}$ 但 $a \neq \{a\}$。
\item \textbf{全集} $U$：在特定讨论范围内包含所有考虑的元素的集合。
\end{itemize}

\subsection{基数}

集合 $A$ 的\textbf{基数}记作 $|A|$。有限集基数是非负整数；无限集中 $\N$ 的基数为 $\aleph_0$（可数无穷），$\R$ 的基数为 $\mathfrak{c} = 2^{\aleph_0}$（不可数）。

\begin{myformula}
\textbf{外延公理}
\[
\forall A \forall B \, (A = B \leftrightarrow \forall x \, (x \in A \leftrightarrow x \in B))
\]
集合由它的元素唯一确定，不依赖于描述方式。
\end{myformula}

\section{子集、真子集、幂集}

\subsection{子集与集合相等}

\begin{mydefinition}{子集}
$A \subseteq B$ 当且仅当 $\forall x (x \in A \rightarrow x \in B)$。
\end{mydefinition}

子集关系的三条基本性质（偏序关系）：
\begin{enumerate}
\item \textbf{自反性}：$A \subseteq A$
\item \textbf{反对称性}：若 $A \subseteq B$ 且 $B \subseteq A$，则 $A = B$
\item \textbf{传递性}：若 $A \subseteq B$ 且 $B \subseteq C$，则 $A \subseteq C$
\end{enumerate}

\begin{mydefinition}{真子集}
$A \subset B$ 当且仅当 $A \subseteq B$ 且 $A \neq B$。
\end{mydefinition}

\textbf{判定集合相等的标准策略}：证明 $A \subseteq B$ 且 $B \subseteq A$，由外延公理得 $A = B$。

\subsection{幂集}

\begin{mydefinition}{幂集}
集合 $A$ 的所有子集构成的集合称为 $A$ 的\textbf{幂集}，记作 $\mathcal{P}(A)$ 或 $2^A$：
\[\mathcal{P}(A) = \{X \mid X \subseteq A\}\]
\end{mydefinition}

\begin{myTheorem}{幂集的基数}
若 $|A| = n$（有限集），则 $|\mathcal{P}(A)| = 2^n$。
\end{myTheorem}

\begin{proof}
\textbf{组合论证}：对 $A$ 中的 $n$ 个元素，每个元素在子集中有"出现"和"不出现"两种选择，因此子集总数为 $2^n$。
\end{proof}

\begin{example}
\begin{itemize}
\item $A = \{1, 2\}$，$\mathcal{P}(A) = \{\varnothing, \{1\}, \{2\}, \{1, 2\}\}$，$|\mathcal{P}(A)| = 4$
\item $A = \varnothing$，$\mathcal{P}(\varnothing) = \{\varnothing\}$，$|\mathcal{P}(\varnothing)| = 1$
\end{itemize}
\end{example}

幂集的性质：
\begin{itemize}
\item $\varnothing \in \mathcal{P}(A)$ 且 $A \in \mathcal{P}(A)$
\item 若 $A \subseteq B$，则 $\mathcal{P}(A) \subseteq \mathcal{P}(B)$
\item $\mathcal{P}(A \cap B) = \mathcal{P}(A) \cap \mathcal{P}(B)$
\end{itemize}

\section{集合运算}

\subsection{并集 $\cup$ 与交集 $\cap$}

\begin{mydefinition}{并集}
$A \cup B = \{x \mid x \in A \lor x \in B\}$
\end{mydefinition}

\begin{mydefinition}{交集}
$A \cap B = \{x \mid x \in A \land x \in B\}$
\end{mydefinition}

若 $A \cap B = \varnothing$，则称 $A$ 和 $B$ \textbf{不相交}。

并集性质：幂等律、交换律、结合律、同一律（$A \cup \varnothing = A$）、零律（$A \cup U = U$）。\\
交集性质：幂等律、交换律、结合律、同一律（$A \cap U = A$）、零律（$A \cap \varnothing = \varnothing$）。

\subsection{差集与补集}

\begin{mydefinition}{差集与补集}
$A - B = A \setminus B = \{x \mid x \in A \land x \notin B\}$。

当给定全集 $U$ 时，$A$ 的\textbf{补集}为 $\overline{A} = U - A = \{x \in U \mid x \notin A\}$。
\end{mydefinition}

补集性质：
\begin{itemize}
\item $\overline{\overline{A}} = A$（双重补）
\item $\overline{U} = \varnothing$，$\overline{\varnothing} = U$
\item $A \cup \overline{A} = U$（排中律）
\item $A \cap \overline{A} = \varnothing$（矛盾律）
\item $\overline{A \cup B} = \overline{A} \cap \overline{B}$（德摩根律）
\item $\overline{A \cap B} = \overline{A} \cup \overline{B}$（德摩根律）
\end{itemize}

\subsection{对称差 $\triangle$}

\begin{mydefinition}{对称差}
$A \triangle B = (A - B) \cup (B - A) = (A \cup B) - (A \cap B)$
\end{mydefinition}

性质：交换律、结合律、$A \triangle \varnothing = A$、$A \triangle A = \varnothing$、$A \triangle B = \varnothing \iff A = B$。

% --- TikZ: 两集合文氏图（并集、交集、差集、对称差）---
\begin{figure}[H]
\centering
\begin{minipage}{0.24\textwidth}
\centering
\begin{tikzpicture}[scale=0.7]
\fill[blue!20] (-1,0) circle (1.5);
\fill[blue!20] (1,0) circle (1.5);
\draw (-1,0) circle (1.5) node[above left=0.8cm] {$A$};
\draw (1,0) circle (1.5) node[above right=0.8cm] {$B$};
\node at (0,-2.2) {并集 $A \cup B$};
\end{tikzpicture}
\end{minipage}\hfill
\begin{minipage}{0.24\textwidth}
\centering
\begin{tikzpicture}[scale=0.7]
\begin{scope}
\clip (-1,0) circle (1.5);
\fill[red!20] (1,0) circle (1.5);
\end{scope}
\draw (-1,0) circle (1.5) node[above left=0.8cm] {$A$};
\draw (1,0) circle (1.5) node[above right=0.8cm] {$B$};
\node at (0,-2.2) {交集 $A \cap B$};
\end{tikzpicture}
\end{minipage}\hfill
\begin{minipage}{0.24\textwidth}
\centering
\begin{tikzpicture}[scale=0.7]
\fill[green!20] (-1,0) circle (1.5);
\begin{scope}
\clip (-1,0) circle (1.5);
\fill[white] (1,0) circle (1.5);
\end{scope}
\draw (-1,0) circle (1.5) node[above left=0.8cm] {$A$};
\draw (1,0) circle (1.5) node[above right=0.8cm] {$B$};
\node at (0,-2.2) {差集 $A - B$};
\end{tikzpicture}
\end{minipage}\hfill
\begin{minipage}{0.24\textwidth}
\centering
\begin{tikzpicture}[scale=0.7]
\fill[orange!20] (-1,0) circle (1.5);
\fill[orange!20] (1,0) circle (1.5);
\begin{scope}
\clip (-1,0) circle (1.5);
\fill[white] (1,0) circle (1.5);
\end{scope}
\draw (-1,0) circle (1.5) node[above left=0.8cm] {$A$};
\draw (1,0) circle (1.5) node[above right=0.8cm] {$B$};
\node at (0,-2.2) {对称差 $A \triangle B$};
\end{tikzpicture}
\end{minipage}
\caption{文氏图——两集合的四种基本运算}
\end{figure}

\subsection{广义并和广义交}

\[
\bigcup_{i=1}^{n} A_i = \{x \mid \exists i, x \in A_i\},\qquad
\bigcap_{i=1}^{n} A_i = \{x \mid \forall i, x \in A_i\}
\]

\section{集合恒等式}

以下列出常用的集合恒等式，每个可用子集法（证明 $左 \subseteq 右$ 且 $右 \subseteq 左$）或命题逻辑等价法证明。

\begin{center}
\begin{tabular}{c|c|c}
\toprule
编号 & 名称 & 恒等式 \\
\midrule
S1 & 幂等律 & $A \cup A = A$；$A \cap A = A$ \\
S2 & 交换律 & $A \cup B = B \cup A$；$A \cap B = B \cap A$ \\
S3 & 结合律 & $(A \cup B) \cup C = A \cup (B \cup C)$ \\
S4 & 分配律 & $A \cup (B \cap C) = (A \cup B) \cap (A \cup C)$ \\
 &  & $A \cap (B \cup C) = (A \cap B) \cup (A \cap C)$ \\
S5 & 同一律 & $A \cup \varnothing = A$；$A \cap U = A$ \\
S6 & 零律 & $A \cup U = U$；$A \cap \varnothing = \varnothing$ \\
S7 & 排中律 & $A \cup \overline{A} = U$ \\
S8 & 矛盾律 & $A \cap \overline{A} = \varnothing$ \\
\bottomrule
\end{tabular}
\end{center}

\begin{myformula}
\textbf{德摩根律}
\[
\overline{A \cup B} = \overline{A} \cap \overline{B},\qquad
\overline{A \cap B} = \overline{A} \cup \overline{B}
\]
\end{myformula}

\begin{center}
\begin{tabular}{c|c|c}
\toprule
编号 & 名称 & 恒等式 \\
\midrule
S10 & 吸收律 & $A \cup (A \cap B) = A$；$A \cap (A \cup B) = A$ \\
S11 & 补余律 & $\overline{\overline{A}} = A$ \\
S12 & 差集恒等式 & $A - B = A \cap \overline{B}$ \\
S13 & 对称差 & $A \triangle B = (A \cup B) - (A \cap B)$ \\
S14 & & $A \cup (B - A) = A \cup B$ \\
S15 & & $A \cap (B - A) = \varnothing$ \\
S16 & & $A - (B \cup C) = (A - B) \cap (A - C)$ \\
S17 & & $A - (B \cap C) = (A - B) \cup (A - C)$ \\
\bottomrule
\end{tabular}
\end{center}

\subsection{恒等式证明示例}

\begin{example}
证明 S4（分配律）：$A \cup (B \cap C) = (A \cup B) \cap (A \cup C)$。

\textbf{（一）证 $\subseteq$}：任取 $x \in A \cup (B \cap C)$。若 $x \in A$，则 $x \in A \cup B$ 且 $x \in A \cup C$，故 $x \in (A \cup B) \cap (A \cup C)$。若 $x \in B \cap C$，则 $x \in B$ 且 $x \in C$，故 $x \in A \cup B$ 且 $x \in A \cup C$。

\textbf{（二）证 $\supseteq$}：任取 $x \in (A \cup B) \cap (A \cup C)$。若 $x \in A$，则 $x \in A \cup (B \cap C)$。若 $x \notin A$，则由 $x \in A \cup B$ 知 $x \in B$；由 $x \in A \cup C$ 知 $x \in C$，故 $x \in B \cap C \subseteq A \cup (B \cap C)$。$\square$
\end{example}

\begin{example}
证明 S9（德摩根律之一）：$\overline{A \cup B} = \overline{A} \cap \overline{B}$。

\[
\begin{aligned}
x \in \overline{A \cup B}
&\iff x \notin A \cup B
\iff \neg(x \in A \lor x \in B) \\
&\iff \neg(x \in A) \land \neg(x \in B)
\iff x \in \overline{A} \land x \in \overline{B} \\
&\iff x \in \overline{A} \cap \overline{B}
\end{aligned}
\]
\end{example}

\section{有序对与笛卡尔积}

\begin{mydefinition}{有序对}
$(a, b) = \{\{a\}, \{a, b\}\}$（库拉托夫斯基定义）。核心性质：
\[(a, b) = (c, d) \iff a = c \land b = d\]
注意 $\{a, b\} = \{b, a\}$ 但 $(a, b) \neq (b, a)$（除非 $a = b$）。
\end{mydefinition}

\begin{mydefinition}{笛卡尔积}
集合 $A$ 与 $B$ 的\textbf{笛卡尔积}为：
\[A \times B = \{(a, b) \mid a \in A \land b \in B\}\]
\end{mydefinition}

$|A \times B| = |A| \cdot |B|$。

笛卡尔积的性质：
\begin{itemize}
\item $A \times \varnothing = \varnothing \times A = \varnothing$
\item 一般 $A \times B \neq B \times A$（不满足交换律）
\item $A \times (B \cup C) = (A \times B) \cup (A \times C)$（分配律）
\item $A \times (B \cap C) = (A \times B) \cap (A \times C)$
\end{itemize}

\subsection{$n$ 元有序组}

递归定义：$(a_1, \dots, a_n) = ((a_1, \dots, a_{n-1}), a_n)$。

\[A_1 \times A_2 \times \dots \times A_n = \{(a_1, \dots, a_n) \mid a_i \in A_i\};\qquad
A^n = \underbrace{A \times A \times \dots \times A}_{n \text{ 个}}\]

例如 $\R^2 = \R \times \R$（平面），$\mathbb{B}^n = \{0,1\}^n$（$2^n$ 个 $n$ 位二进制串）。

% --- TikZ: 笛卡尔积示意图 ---
\begin{figure}[H]
\centering
\begin{tikzpicture}[scale=1.2]
\draw[-{Stealth}] (-0.5,0) -- (5.5,0) node[right] {$A$};
\draw[-{Stealth}] (0,-0.5) -- (0,4.5) node[above] {$B$};
% A = {1,2,3,4,5}, B = {1,2,3,4}
\foreach \x in {1,...,5} {
    \draw (\x,0.05) -- (\x,-0.05) node[below] {\footnotesize $\x$};
    \foreach \y in {1,...,4} {
        \fill[blue!60] (\x,\y) circle (1.5pt);
    }
}
\foreach \y in {1,...,4} {
    \draw (0.05,\y) -- (-0.05,\y) node[left] {\footnotesize $\y$};
}
\node[font=\small] at (3,2.5) {$(a,b) \in A \times B$};
\node[font=\small] at (5.2,-0.3) {$|A|=5$};
\node[font=\small] at (-0.3,4.2) {$|B|=4$};
\node[font=\small] at (5.2,4) {$|A \times B| = 20$};
\end{tikzpicture}
\caption{笛卡尔积 $A \times B$ 的几何表示——每个点对应一个有序对}
\end{figure}

\section{容斥原理}

\subsection{两集合与三集合}

\begin{myTheorem}{两集合容斥原理}
\[|A \cup B| = |A| + |B| - |A \cap B|\]
\end{myTheorem}

直观解释：分别计算 $|A|$ 和 $|B|$ 时，$A \cap B$ 被计算了两次，因此需要减去一次。

\begin{myTheorem}{三集合容斥原理}
\[
\begin{aligned}
|A \cup B \cup C| &= |A| + |B| + |C| \\
&\quad - |A \cap B| - |A \cap C| - |B \cap C| \\
&\quad + |A \cap B \cap C|
\end{aligned}
\]
\end{myTheorem}

% --- TikZ: 三集合容斥原理文氏图 ---
\begin{figure}[H]
\centering
\begin{tikzpicture}[scale=1.1]
\tikzset{
    dmVenn3/.style={circle,draw=blue!60,fill=blue!10,minimum size=3.5cm,opacity=0.5}
}
\node[dmVenn3] (A) at (0,1.2) {};
\node[dmVenn3] (B) at (-1.2,-0.5) {};
\node[dmVenn3] (C) at (1.2,-0.5) {};
\node at (0,2.0) {$A$};
\node at (-1.8,-0.8) {$B$};
\node at (1.8,-0.8) {$C$};
\node[font=\small] at (0,0.3) {$A \cap B \cap C$};
\node[font=\small] at (0,-1.3) {$|A \cup B \cup C| = \text{加单}- \text{减双}+ \text{加三}$};
\end{tikzpicture}
\caption{三集合容斥原理文氏图——先加单集，减去两两交集，再加回三集交集}
\end{figure}

\subsection{一般 $n$ 个集合的容斥原理}

\begin{myformula}
\textbf{容斥原理一般形式}
\[
\left|\bigcup_{i=1}^{n} A_i\right| = \sum_{k=1}^{n} (-1)^{k-1} \sum_{1 \leq i_1 < \dots < i_k \leq n} |A_{i_1} \cap A_{i_2} \cap \dots \cap A_{i_k}|
\]
符号规则：奇数个集合的交集前符号为 $+$，偶数个为 $-$。
\end{myformula}

\textbf{补集形式}：
\[
\left|\bigcap_{i=1}^{n} \overline{A_i}\right| = |U| - \left|\bigcup_{i=1}^{n} A_i\right|
\]

\subsection{应用示例}

\begin{example}
1 到 100 中能被 3 或 5 整除的整数有多少个？

$|A| = \lfloor 100/3 \rfloor = 33$，$|B| = \lfloor 100/5 \rfloor = 20$，
$|A \cap B| = \lfloor 100/15 \rfloor = 6$。

$|A \cup B| = 33 + 20 - 6 = 47$。
\end{example}

\begin{example}
三集合问题：某班 50 人，学英语 30 人，学日语 20 人，学法语 15 人；学英+日 10 人，学英+法 7 人，学日+法 5 人；三门都学 3 人。问至少学一门的有多少人？

$|E \cup J \cup F| = 30 + 20 + 15 - 10 - 7 - 5 + 3 = 46$ 人。
\end{example}

% --- TikZ: 幂集哈塞图 ---
\begin{figure}[H]
\centering
\begin{tikzpicture}[node distance=1.2cm, auto]
\tikzset{
    dmHasse/.style={circle,draw=blue!60,fill=blue!10,minimum size=6mm,font=\small}
}
\node[dmHasse] (123) at (0,0) {$\{a,b,c\}$};
\node[dmHasse] (12) at (-1.5,1.5) {$\{a,b\}$};
\node[dmHasse] (13) at (0,1.5) {$\{a,c\}$};
\node[dmHasse] (23) at (1.5,1.5) {$\{b,c\}$};
\node[dmHasse] (1) at (-1.5,3) {$\{a\}$};
\node[dmHasse] (2) at (0,3) {$\{b\}$};
\node[dmHasse] (3) at (1.5,3) {$\{c\}$};
\node[dmHasse] (0) at (0,4.5) {$\varnothing$};
\draw[thick] (123) -- (12);
\draw[thick] (123) -- (13);
\draw[thick] (123) -- (23);
\draw[thick] (12) -- (1);
\draw[thick] (12) -- (2);
\draw[thick] (13) -- (1);
\draw[thick] (13) -- (3);
\draw[thick] (23) -- (2);
\draw[thick] (23) -- (3);
\draw[thick] (1) -- (0);
\draw[thick] (2) -- (0);
\draw[thick] (3) -- (0);
\node[font=\footnotesize] at (0,5.3) {$\mathcal{P}(\{a,b,c\})$ 的哈塞图（按 $\subseteq$ 偏序）};
\end{tikzpicture}
\caption{幂集 $\mathcal{P}(\{a,b,c\})$ 的哈塞图——共 $2^3=8$ 个子集}
\end{figure}

\begin{mycode}{Python集合运算与容斥原理}
\begin{lstlisting}[language=Python]
from itertools import combinations
from typing import Set, List, Any

A = {1, 2, 3, 4, 5}
B = {4, 5, 6, 7, 8}

# 基本运算
union = A | B              # {1,2,3,4,5,6,7,8}
inter = A & B              # {4,5}
diff  = A - B              # {1,2,3}
sym   = A ^ B              # {1,2,3,6,7,8}
print(f"A ⊆ B: {A <= B}")  # False

def cartesian_product(*sets: Set[Any]) -> Set[tuple]:
    result = {(x,) for x in sets[0]}
    for s in sets[1:]:
        result = {r + (y,) for r in result for y in s}
    return result

def powerset(s: Set[Any]) -> Set[frozenset]:
    elements = list(s)
    n = len(elements)
    return {frozenset(elements[i] for i in range(n)
            if (mask >> i) & 1) for mask in range(1 << n)}

def inclusion_exclusion(sets: List[Set[Any]]) -> int:
    n = len(sets)
    total = 0
    for k in range(1, n + 1):
        sign = (-1) ** (k - 1)
        for indices in combinations(range(n), k):
            inter = sets[indices[0]]
            for idx in indices[1:]:
                inter = inter & sets[idx]
            total += sign * len(inter)
    return total

# 验证容斥原理
div3 = {x for x in range(1, 101) if x % 3 == 0}
div5 = {x for x in range(1, 101) if x % 5 == 0}
print(f"|3的倍数 ∪ 5的倍数| = {inclusion_exclusion([div3, div5])}")  # 47
\end{lstlisting}
\end{mycode}
